Las instrucciones de uso del \textit{software} se detallan a continuación. Este manual se dirige al usuario final del \textit{software} objetivo de este proyecto.

\section{Introducción}

MoneyBall LaLiga es una aplicación móvil orientada a la consulta y análisis de información deportiva de LaLiga. El sistema permite explorar temporadas, equipos, jugadores y partidos mediante clasificaciones, estadísticas, gráficos y diferentes herramientas de análisis.

La aplicación está dirigida principalmente a aficionados al fútbol, analistas deportivos y usuarios interesados en estudiar el rendimiento de los participantes de la competición. Además, incorpora funciones como la gestión de favoritos, un asistente conversacional y un simulador de resultados, sin requerir conocimientos técnicos o estadísticos avanzados.

\section{Uso del sistema}

En esta sección se describen las principales acciones que puede realizar un usuario dentro de la aplicación. La mayoría de las pantallas comparten una estructura similar, formada por una cabecera, un menú lateral, listas de elementos y pestañas que permiten acceder a distintos tipos de información.

\subsection{Registro e inicio de sesión}

Al abrir la aplicación, el usuario accede a la pantalla de inicio de sesión. Para entrar en el sistema debe introducir su nombre de usuario o correo electrónico y su contraseña, y pulsar el botón \texttt{Iniciar sesión}. El icono situado junto al campo de contraseña permite mostrar u ocultar los caracteres introducidos.

Si el usuario todavía no dispone de una cuenta, puede seleccionar la opción \texttt{Regístrate aquí}. En el formulario de registro deberá indicar un nombre de usuario, un correo electrónico y una contraseña, además de confirmar esta última. La contraseña debe contener al menos ocho caracteres, una letra mayúscula, una letra minúscula y un número.

Una vez completado correctamente el registro, el usuario puede volver a la pantalla de acceso e iniciar sesión con sus credenciales. Si algún dato es incorrecto o no se puede establecer conexión con el servidor, la aplicación muestra un mensaje informativo.

\subsection{Navegación por la aplicación}

Después de iniciar sesión se muestra la pantalla de inicio. El botón situado en la parte superior izquierda abre el menú lateral, desde el cual se puede acceder a las siguientes secciones:

\begin{itemize}
    \item Inicio: muestra un resumen de la información más relevante para el usuario.
    \item Análisis IA: permite acceder a los análisis avanzados y las predicciones.
    \item Temporadas: contiene la información histórica de las temporadas disponibles.
    \item Equipos: permite consultar los clubes registrados en el sistema.
    \item Jugadores: permite localizar y analizar futbolistas.
    \item Partidos: ofrece la consulta y comparación de encuentros.
    \item Simulador: permite crear escenarios alternativos para la temporada actual.
\end{itemize}

La cabecera también incluye un botón de búsqueda. Al pulsarlo, el usuario puede introducir el nombre de un equipo o jugador y acceder directamente a su pantalla de detalle. Para llevar a cabo la búsqueda es necesario escribir al menos dos caracteres.

En las pantallas que contienen varias categorías de información se utilizan pestañas superiores. El usuario puede cambiar de pestaña pulsando sobre su nombre o desplazándose horizontalmente.

\subsection{Consulta y análisis de información deportiva}

La aplicación organiza la información deportiva en cuatro ámbitos principales: temporadas, equipos, jugadores y partidos. El procedimiento general de consulta es similar en todos ellos: el usuario accede a una sección desde el menú, localiza el elemento deseado y pulsa sobre él para abrir su pantalla de detalle.

En la sección de temporadas se pueden consultar los partidos disputados, la clasificación, los equipos participantes, los rankings y los gráficos generales de cada edición de la competición.

La sección de equipos permite buscar un club y acceder a información sobre sus partidos, trayectoria, clasificación, plantilla y rendimiento. Los paneles de análisis muestran indicadores y gráficos que ayudan a interpretar la evolución y las características de juego del equipo.

En la sección de jugadores se puede efectuar una búsqueda por nombre y consultar información general, estadísticas, atributos, partidos y trayectoria deportiva. Los filtros disponibles permiten reducir los resultados y facilitar la comparación entre futbolistas.

La sección de partidos permite seleccionar dos equipos para consultar su historial de enfrentamientos. Al acceder al detalle de un encuentro se muestran diferentes apartados, como la previa, las alineaciones, los eventos y el análisis posterior al partido.

Los gráficos, tablas y tarjetas estadísticas pueden requerir desplazamiento vertical u horizontal para visualizar todo su contenido. Cuando una pantalla ofrece filtros, el usuario puede modificar sus valores para actualizar la información representada.

\subsection{Gestión de favoritos}

Los equipos y jugadores pueden marcarse como favoritos mediante el botón con forma de corazón que aparece en los listados y en las pantallas de detalle. Al pulsarlo, el elemento se añade a la selección personal del usuario y el icono cambia para indicar su nuevo estado.

Para eliminar un favorito basta con pulsar nuevamente el mismo botón. Los cambios quedan asociados a la cuenta y se utilizan para personalizar determinadas secciones, como la pantalla de inicio y el apartado de análisis de favoritos.

\subsection{Uso del apartado Análisis IA}

La sección Análisis IA reúne diferentes herramientas para interpretar los datos almacenados en el sistema. Antes de realizar una consulta, el usuario puede seleccionar la temporada que desea analizar.

El contenido se organiza en las siguientes pestañas:

\begin{itemize}
    \item Favoritos: presenta información conjunta sobre los equipos y jugadores que el usuario ha marcado como favoritos.
    \item Equipos: permite seleccionar equipos y analizar su rendimiento mediante indicadores, comparaciones y predicciones.
    \item Jugadores: permite buscar y comparar futbolistas aplicando diferentes filtros deportivos.
    \item Predicciones: muestra los resultados generados por los modelos de análisis y minería de datos.
\end{itemize}

Para utilizar estas funciones, el usuario debe seleccionar la temporada y completar los filtros solicitados en cada pestaña. Los resultados se muestran mediante tarjetas, tablas y gráficos. Estos análisis sirven como apoyo para interpretar los datos y no modifican la información oficial almacenada en el sistema.

\subsection{Uso del asistente Agente de Oro}

El \textit{Agente de Oro} es el asistente conversacional de la aplicación. Puede abrirse pulsando el botón flotante dorado que aparece en la esquina inferior de las principales pantallas.

Una vez abierto, el usuario puede escribir una pregunta relacionada con equipos, jugadores, partidos, clasificaciones o rankings. También puede seleccionar una de las sugerencias mostradas por el asistente. Al pulsar el botón de envío, la consulta se procesa y la respuesta aparece dentro de la conversación.

El asistente conserva el contexto de los últimos mensajes, por lo que es posible realizar preguntas relacionadas con una consulta anterior. La opción \texttt{Nuevo chat} elimina el contexto de la conversación actual y permite comenzar una consulta independiente.

Las respuestas se generan a partir de los datos disponibles en la aplicación. Si la pregunta es ambigua, no existen datos suficientes o se produce un problema de conexión, el asistente puede mostrar un mensaje de error o solicitar una consulta más concreta.

\subsection{Uso del simulador manual}

El simulador permite estudiar cómo determinados resultados podrían afectar a la clasificación de la temporada actual. Los cambios realizados son temporales y no modifican los resultados oficiales.

La pantalla se divide en las pestañas \texttt{Resultados} y \texttt{Clasificación}. En la primera, el usuario puede desplazarse entre las jornadas e introducir marcadores para los partidos pendientes. Los encuentros que ya han finalizado aparecen bloqueados y no pueden modificarse.

Después de introducir uno o varios marcadores completos, se debe pulsar \texttt{Recalcular escenario}. La aplicación actualizará la clasificación teniendo en cuenta los resultados simulados. El botón de restablecimiento permite eliminar los marcadores introducidos y recuperar el escenario inicial.

En la pestaña \texttt{Clasificación} se puede consultar la tabla resultante. Dentro de esta pantalla, la opción \texttt{Probabilidades} permite ejecutar una simulación que calcula las posibilidades de cada equipo de ser campeón, clasificarse para competiciones europeas, finalizar en la zona media o descender.

Los resultados de la simulación se mantienen únicamente durante el uso del simulador y no alteran los datos reales de la competición.

\subsection{Gestión del perfil de usuario}

El perfil se abre pulsando sobre el nombre o el icono del usuario situado en la parte superior del menú lateral. Esta pantalla muestra el nombre de usuario y el correo electrónico asociados a la cuenta.

La opción \texttt{Editar perfil} permite modificar estos datos. Una vez efectuados los cambios, el usuario debe pulsar \texttt{Guardar cambios} para actualizar la información.

Desde el perfil también se puede acceder a \texttt{Cambiar contraseña}. Para completar esta operación es necesario introducir la contraseña actual, la nueva contraseña y su confirmación. La nueva contraseña debe cumplir los requisitos de seguridad indicados por la aplicación.

También se puede eliminar la cuenta pulsando el botón de \texttt{Eliminar cuenta}, una vez pulsado, el usuario debe confirmar la operación asumiendo las consecuencias que le trae.

\subsection{Ajustes y cierre de sesión}

La sección \texttt{Ajustes}, disponible en el menú lateral, permite activar o desactivar el modo oscuro. La preferencia seleccionada se conserva aunque el usuario cierre la aplicación.

Para finalizar la sesión, se puede utilizar el botón \texttt{Cerrar sesión} del menú lateral o la opción equivalente disponible en el perfil. Al cerrar la sesión, se eliminan los datos de acceso almacenados localmente y se vuelve a mostrar la pantalla de inicio de sesión.



